
\section{Literature Review}

This study sits at the intersection of empirical AI-deception research, financial-domain LLM evaluation, and agent-safety monitoring. Its central claim is that deception in financial LLM agents should not be studied only as explicit lying under contrived prompts, but as a deployment-relevant failure mode in which an agent facing realistic goal conflict conceals, downplays, or misframes material risk. The work therefore prioritises low-nudge financial goal-conflict settings, treats risk concealment as the central deception measure, and makes detectability a primary contribution rather than an afterthought.

Financial services are a strong test domain because LLM outputs are not merely informational; they can influence investment recommendations, compliance judgements, customer communications, trading workflows, and risk escalation. UK regulators report substantial AI adoption in financial services, including use cases involving automated decision-making and partial firm understanding of deployed AI technologies \citep{bankofengland2024ai}. The FCA's current AI approach similarly emphasises safe adoption, governance, accountability, and consumer outcomes \citep{fca2025ai}. This makes finance an unusually useful domain for deception research: the harms of omission or misleading framing are concrete, but the relevant behaviours may be too subtle for generic truthfulness or jailbreak benchmarks.

The finance-LLM literature initially focused on capability and domain adaptation. BloombergGPT showed that a large model trained on financial and general corpora could perform strongly on financial tasks while retaining general capabilities \citep{wu2023bloomberggpt}. FinGPT responded from the open-source side, arguing that financial LLM development depends on data-centric pipelines, real-time financial data, and accessible adaptation methods \citep{yang2023fingpt}. FinBen broadened evaluation across information extraction, textual analysis, question answering, text generation, risk management, forecasting, and decision-making, while finding that models still struggle with advanced reasoning and complex financial tasks \citep{xie2024finben}. More recent work on LLM agents in finance shifts the unit of analysis from static model answers to deployed workflows across data analysis, investment research, trading, investment management, and risk management \citep{li2025financialagents}. The trajectory is clear: finance research has moved from asking whether LLMs understand financial text to asking whether agentic systems can act safely in financial institutions.

A key conceptual distinction is between falsehood, error, sycophancy, and deception. TruthfulQA established that language models may produce false answers learned from human misconceptions, but such behaviour is not necessarily strategic; it can arise from imitation of flawed text distributions \citep{lin2022truthfulqa}. Work on sycophancy similarly shows that preference optimisation can encourage models to match user beliefs over truthful answers, because human feedback may reward agreeable responses \citep{sharma2023sycophancy}. These are adjacent failure modes, but they are not sufficient for the current project. A finance agent that fails to mention downside risk may be mistaken, overly agreeable, or strategically concealing adverse information. The methodological challenge is therefore to design scenarios where the model has access to risk-relevant facts, the stakeholder would reasonably expect disclosure, and the final output can be compared against known evidence.

The broader AI-deception literature provides useful definitions and warnings. \citet{park2023deception} define deception as the systematic inducement of false beliefs in pursuit of an outcome other than truth, and argue that both special-purpose systems and general-purpose systems can display deceptive behaviour. Learned-optimisation work supplies the older theoretical foundation: if a learned system becomes an optimiser with an objective different from the training objective, it may appear aligned during evaluation while optimising differently under deployment conditions \citep{hubinger2019risks}. However, the current study should avoid overclaiming that financial LLM agents are mesa-optimisers. The more defensible empirical claim is that, under goal conflict, current LLM agents may produce outputs that induce false beliefs about financial risk, even without being explicitly told to lie.

Empirical deception research has already shown that model behaviour can become strategically misleading under pressure. \citet{scheurer2023strategically} deploy GPT-4 as an autonomous trading agent that receives an insider tip, acts on it despite knowing management disapproves, and then hides the genuine reason for the trade when reporting to a manager. This paper is highly relevant because it is already financial in flavour and because the deception is not simply a response to a user instruction to lie. Yet it remains high-pressure and stylised: the model faces an insider-trading dilemma where deception is salient and the experimental design intentionally creates acute pressure. This study can improve ecological validity by using lower-nudge scenarios where the conflict is more mundane: preserving a favourable recommendation, meeting a business objective, avoiding escalation, or maintaining a clean audit trail while adverse risk facts are available.

Recent scheming and agentic-misalignment studies sharpen the concern. Work on in-context scheming evaluates models in agentic settings where they are given goals and environments that incentivise covert pursuit, reporting behaviours such as subtle mistakes, attempts to disable oversight, and persistence in deception during follow-up interrogation \citep{meinke2024scheming}. Sleeper Agents shows that deliberately trained deceptive behaviours can persist through standard safety training, including supervised fine-tuning, reinforcement learning, and adversarial training \citep{hubinger2024sleeper}. Agentic Misalignment stress-tests models in hypothetical corporate environments with email access and sensitive information, finding that models can resort to insider-like behaviours such as blackmail or leaking information when facing replacement or goal conflict \citep{anthropic2025agenticmisalignment}. Alignment Faking reveals a complementary pattern: when presented with a system prompt indicating that training was active, Claude~3 Opus selectively complied with objectionable queries it would otherwise refuse, apparently to avoid modification of its values, with this strategic compliance disappearing outside that perceived-training context \citep{greenblatt2024alignmentfaking}. This behaviour is distinct from hidden-goal scheming; it arises from the model's inference about evaluation pressure, suggesting that models may adapt their conduct to the perceived stakes of the interaction. These papers justify taking agentic deception seriously, but they also highlight a gap: many evaluations rely on extreme stakes, explicit hidden goals, or dramatic misconduct. Financial risk concealment may be more subtle and more common precisely because it resembles ordinary professional communication.

This is where low-nudge deception becomes central to the present study. \citet{wu2025beyond} provide the primary methodological anchor for this project's experimental philosophy. They criticise the tendency of prior work to induce deception through hidden-objective prompting or fine-tuning, and instead study self-initiated deception on benign prompts. Their framework separates deceptive intention from deceptive behaviour and shows that both can rise with task difficulty, establishing that deception is not an artefact of adversarial elicitation but can emerge under ordinary interaction conditions. This precedent is not merely motivational; it directly grounds the present study's low-nudge design, in which goal conflict arises from the financial task itself rather than from an explicit instruction to deceive. A financial agent does not need to be told to hide a regulatory issue for risk concealment to matter; it may be enough that the system is asked to produce a favourable client-facing summary, justify a recommendation, or avoid unnecessary escalation.

Finance-specific safety work increasingly supports a risk-first framing. \citet{chen2025riskfirst} argue that standard financial LLM-agent benchmarks prioritise task performance while neglecting safety risks such as hallucination, stale data, and adversarial vulnerabilities. FinVault introduces execution-grounded financial-agent safety scenarios with state-writable databases, compliance constraints, regulatory cases, real-world vulnerabilities, and benign inputs for false-positive evaluation \citep{finvault2026}. FinSafetyBench evaluates financial safety using real-world crime cases and ethics standards across financial-crime and ethical-violation categories \citep{finsafetybench2026}. FinMCP-Bench evaluates financial tool use through Model Context Protocol tools, covering single-tool, multi-tool, and multi-turn tasks across several financial scenarios \citep{finmcp2026}. Together, these benchmarks show that finance-agent evaluation is becoming more realistic, tool-integrated, and safety-aware. Their limitation for the current research question is that they mostly evaluate harmful requests, tool-use accuracy, prompt injection, regulatory compliance, or state-changing security failures. They do not fully isolate whether a model with access to adverse facts will voluntarily disclose those facts when disclosure conflicts with an assigned financial objective.

The closest direct anchor is \citet{cheng2025riskconcealment}, who introduce Risk-Concealment Attacks, a multi-turn framework that iteratively hides regulatory risks in order to elicit seemingly compliant but regulatory-violating responses. This paper is important because it treats risk concealment as a finance-native safety failure. However, it is adversarial red-teaming: the user or attack framework actively conceals risk from the model. The study asks the complementary question: when the model itself has access to material risk facts, does it conceal them from a stakeholder under low-nudge goal conflict? This is an important distinction. Risk-concealment attacks study susceptibility to external manipulation; low-nudge financial deception studies whether the agent becomes the source of concealment.

General agent-safety benchmarks also inform the experimental design. ToolEmu demonstrates a scalable way to test language-model agents with emulated tools and risk evaluators \citep{ruan2023toolemu}. AgentDojo shows that agents executing tools over untrusted data are vulnerable to prompt injection and that realistic task environments can challenge both models and defences \citep{debenedetti2024agentdojo}. AgentHarm and Agent-SafetyBench broaden the safety perspective by evaluating multi-step harmful agent tasks and diverse interaction environments \citep{andriushchenko2024agentharm,zhang2024agentsafetybench}. These benchmarks justify using scenario-based, multi-step, workflow-aware evaluations rather than static question answering. Yet they leave room for this dissertation's contribution: subtle stakeholder-facing deception in legitimate financial work is not the same as malicious-user compliance, prompt injection, or overt harmful task completion.

The taxonomy should therefore foreground omissions and framing rather than direct lies. Recent work on auditing LLM deception through information manipulation theory formalises this intuition, operationalising deception as the manipulation of atomic information units---including selective omission, distortion of quantity or relevance, and misleading framing---rather than as outright falsehood \citep{decor2026}. In finance, materiality often depends on what a reasonable customer, manager, auditor, or compliance reviewer would expect to know. An output can be deceptive while remaining technically true: it may present upside scenarios while excluding downside stress tests; mention ``manageable regulatory considerations'' while omitting a known enforcement risk; claim that ``standard checks were completed'' when the record shows no escalation; or answer follow-up questions by rationalising the original recommendation rather than acknowledging omitted evidence. This focus is consistent with the hypothesis that risk concealment may be more common than direct lying. It also better reflects financial professional norms: deception often occurs through selective disclosure, not easily falsifiable statements.

Detection is the second major contribution. \citet{pacchiardi2023catch} propose a black-box lie detector that asks unrelated follow-up questions after a suspected lie and uses yes/no answers in a classifier. Their finding that such signals can generalise across settings supports the idea that deceptive outputs may leave behavioural traces beyond the original statement. However, subsequent work probing the limits of this lie-detector paradigm demonstrates that its effectiveness degrades when deception operates through true-but-misleading outputs, selective omission, or framing effects rather than through falsifiable assertions \citep{probinglimits2026}. If an agent conceals a material risk by simply not mentioning it, or buries it under favourable framing, there is no discrete lie for a follow-up classifier to detect. The information-manipulation framing of DECOR \citep{decor2026}, which decomposes outputs into atomic information units and audits whether relevant units have been omitted or distorted, suggests that detection must move beyond binary lie classification toward structured comparison of what the model had access to and what it chose to disclose. At the same time, finance requires more than generic lie detection. An auditor needs to know whether the recommendation contradicts known facts, whether adverse evidence was omitted, whether disclosures are balanced, and whether follow-up explanations remain consistent with the record. 

% proposed detector TBC: The proposed audit-style checks---comparing known risk facts, final recommendations, disclosures, and post-hoc explanations---are therefore more deployment-relevant than a purely generic detector.

White-box detection offers a useful comparator. \citet{goldowskydill2025linearprobes} train linear probes on model activations to detect strategic deception, including insider-trading concealment and safety-evaluation sandbagging. \citet{natarajan2026deceptionprobes} then show that probe performance depends heavily on the training instruction pair and that targeted deception taxonomies improve detection. This argues against a universal deception detector and in favour of threat-model-specific probes. 

% proposed detector TBC: For this study, a finance-specific taxonomy centred on risk concealment could generate better honest/deceptive contrast pairs than generic truthfulness prompts. However, many deployed financial LLM systems will be API-based and may not expose activations, so black-box behavioural detection and audit-log consistency checks should remain primary.

A further validity concern for any deception evaluation arises from work on AI sandbagging, which demonstrates that language models can be prompted or fine-tuned to strategically underperform on capability evaluations while maintaining normal performance in other contexts \citep{vanderweij2024sandbagging}. The implication for deception research is direct: if models can modulate their behaviour based on whether they detect an evaluation setup, they may also produce more honest outputs under perceived audit scrutiny than they would in routine deployment. Combined with the evaluation-pressure sensitivity revealed by alignment faking \citep{greenblatt2024alignmentfaking}, this concern reinforces the importance of low-nudge, deployment-realistic scenarios---where the evaluation context is not telegraphed---and argues against relying solely on test-time behavioural signals that a model could learn to suppress.

Interpretability and latent-knowledge work further clarifies why detection should compare what the model had access to with what it says. \citet{burns2022discovering} attempt to recover latent knowledge from model activations without supervision, motivated by the gap between what models know internally and what they output. Inference-Time Intervention similarly suggests that models may encode truth-relevant directions that can be shifted to improve truthful answers, though with trade-offs between truthfulness and helpfulness \citep{li2023inference}. Representation engineering extends this logic to high-level traits such as honesty, hallucination, sycophancy, and harmfulness \citep{zou2023representation}. These methods are promising, but the dissertation should treat them as complementary unless activation access is available. For a finance deployment study, the strongest first-line evidence will come from observable inconsistencies: the model saw a material risk, produced a recommendation or disclosure that concealed it, and then defended or minimised the omission when questioned.

% Mitigation is sensibly framed as a stretch goal because the literature suggests simple instruction-based fixes may be fragile. Prior work varies system instructions, pressure, scratchpad access, and perceived detection risk, indicating that small environmental changes can affect but not eliminate deceptive behaviour \citep{scheurer2023strategically}. Sleeper Agents warns that safety training can fail to remove hidden deceptive behaviours and may even teach models to better recognise triggers \citep{hubinger2024sleeper}. Governance-oriented controls may therefore be more realistic for this project than claims about fully eliminating deception. Escalation-channel work is especially relevant because it treats escalation as an environmental control for task-rule conflict, evaluating whether agents choose safer actions when provided with structured routes to surface conflicts \citep{gomez2025escalation}. In finance, analogous controls include mandatory material-risk disclosure, uncertainty statements, audit-evidence references, escalation options, and incentives for saying that a recommendation cannot be made without compliance review.

% TBC: Overall, the literature justifies the study on three grounds. First, existing deception research shows that LLM agents can behave deceptively under pressure, in scheming environments, and in benign-prompt settings, but much of the evidence is high-nudge, game-like, adversarial, or not finance-specific. Second, financial LLM research has matured from capability benchmarking to agentic, execution-grounded, and risk-aware evaluation, but still under-measures voluntary concealment of known material risks. Third, detection research shows promise in both black-box behavioural classifiers and white-box probes, yet finance needs audit-style detection grounded in factual records, disclosures, and follow-up consistency. The proposed study's contribution is therefore not merely another deception benchmark. It reframes deception as a financial governance problem: when an LLM agent has access to material adverse evidence and faces ordinary goal conflict, does it disclose the risk, conceal it, or misframe it, and can that behaviour be detected before it affects a stakeholder?

%% ============================================================
%% NEW BIB ENTRIES — append to references.bib
%% ============================================================
%
% @article{decor2026,
%   title        = {DECOR: Auditing LLM Deception via Information Manipulation Theory},
%   author       = {{DECOR Authors}},
%   journal      = {arXiv preprint arXiv:2605.19270},
%   year         = {2026},
%   url          = {https://arxiv.org/abs/2605.19270}
% }
%
% @article{probinglimits2026,
%   title        = {Probing the Limits of the Lie Detector Approach to LLM Deception},
%   author       = {{Probing Limits Authors}},
%   journal      = {arXiv preprint arXiv:2603.10003},
%   year         = {2026},
%   url          = {https://arxiv.org/abs/2603.10003}
% }
%
% @article{greenblatt2024alignmentfaking,
%   title        = {Alignment Faking in Large Language Models},
%   author       = {Greenblatt, Ryan and others},
%   journal      = {arXiv preprint arXiv:2412.14093},
%   year         = {2024},
%   url          = {https://arxiv.org/abs/2412.14093}
% }
%
% @article{vanderweij2024sandbagging,
%   title        = {AI Sandbagging: Language Models Can Strategically Underperform on Evaluations},
%   author       = {van der Weij, Teun and others},
%   journal      = {arXiv preprint arXiv:2406.07358},
%   year         = {2024},
%   url          = {https://arxiv.org/abs/2406.07358}
% }
%% ============================================================
